\begin{solapartat}
Si el protó no es desvia, la força resultant és zero, $F_e = F_m$

\punts{0,4} Esquema

\figura[0.55]{sol_fig1.png}

$B = 3,00\times 10^{3}\ \mathrm{G}\times\dfrac{1\ \mathrm{T}}{10^{4}\ \mathrm{G}} = 3,00\times 10^{-1}\ \mathrm{T}$

\punts{0,2} Perquè el protó no es desviï de la seva trajectòria: $F_e = F_m$

\punts{0,1} $\left.\begin{aligned} F_e &= qE\\ F_m &= qvB\end{aligned}\right\}\ v = \dfrac{E}{B}$

\punts{0,3} $v = \dfrac{2,00\times 10^{5}}{3,00\times 10^{-1}} = 6,67\times 10^{5}\ \mathrm{m/s} = 667\ \mathrm{km/s}$
\end{solapartat}

\begin{solapartat}
Si només hi ha camp magnètic, el protó es desviarà seguint una trajectòria circular per l'efecte de la força magnètica.

\punts{0,2} $\vec{F}_m = m\vec{a}$

\punts{0,2} $qvB = m\,\dfrac{v^2}{R} \Rightarrow R = \dfrac{mv}{qB}$

% DUBTE: amb aquestes dades el resultat és 2,32×10^-2 m (2,32 cm); la pauta diu 2,30.
\punts{0,6} $R = \dfrac{(1,67\times 10^{-27})\times(6,67\times 10^{5})}{(1,60\times 10^{-19})\times(3,00\times 10^{-1})} = 2,30\times 10^{-2}\ \mathrm{m} = 2,30\ \mathrm{cm}$
\end{solapartat}
